\documentclass[11pt,a4paper]{article}

% ---------------------------------------------------------------------------
% Codificación e idioma
% ---------------------------------------------------------------------------
\usepackage[utf8]{inputenc}
\usepackage[T1]{fontenc}
\usepackage[spanish,es-nodecimaldot]{babel}

% ---------------------------------------------------------------------------
% Matemática, tipografía y utilidades
% ---------------------------------------------------------------------------
\usepackage{amsmath,amssymb}
\usepackage{booktabs}
\usepackage{array}
\usepackage{enumitem}
\usepackage{geometry}
\geometry{margin=2.5cm}
\usepackage{xcolor}
\usepackage[hidelinks]{hyperref}
\hypersetup{
  colorlinks=true,
  linkcolor=blue!50!black,
  citecolor=blue!50!black,
  urlcolor=blue!60!black,
  pdftitle={QTransLIFIA: traducción y verificación semántica de circuitos cuánticos desde Quirk a ecosistemas de software cuántico},
  pdfkeywords={Quirk; OpenQASM; ingeniería de software cuántica; preservación semántica; cálculo ZX}
}

% ---------------------------------------------------------------------------
% Metadatos del documento
% ---------------------------------------------------------------------------
\title{\bfseries QTransLIFIA: traducción y verificación semántica de circuitos \\
cuánticos desde Quirk hacia ecosistemas de software cuántico}

\author{
  Equipo del Proyecto QTransLIFIA\\
  \small Laboratorio de Investigación y Formación en Informática Avanzada (LIFIA)
}
\date{\today}

% Comandos de conveniencia para código y términos técnicos
\newcommand{\code}[1]{\texttt{#1}}
\newcommand{\gate}[1]{\textit{#1}}

\begin{document}
\maketitle

% ---------------------------------------------------------------------------
% Resumen y palabras clave
% ---------------------------------------------------------------------------
\begin{abstract}
\noindent
La computación cuántica se apoya en un ecosistema fragmentado de herramientas de diseño,
lenguajes intermedios y \emph{software development kits} (SDK) de ejecución que no
comparten una representación común. Esta heterogeneidad dificulta la portabilidad de los
circuitos y amenaza la fidelidad semántica del diseño original cuando este debe migrarse
entre plataformas. En este trabajo presentamos \textbf{QTransLIFIA}, un traductor de
circuitos cuánticos que convierte diseños elaborados en la herramienta interactiva
\emph{Quirk} hacia \emph{OpenQASM} como lenguaje pivote de representación intermedia y,
desde allí, hacia múltiples ecosistemas de destino (IBM Qiskit, AWS Braket y Xanadu
PennyLane). A diferencia de una transliteración sintáctica, QTransLIFIA adopta un enfoque
de \emph{traducción semántica con preservación de abstracciones}: conserva las compuertas
personalizadas definidas por el usuario ---tanto las compuestas por circuito como las
definidas por matriz---, la disposición columnar y las relaciones de control y objetivo del
circuito original. El trabajo se enmarca en la \emph{Ingeniería de Software Cuántica}
(\emph{Quantum Software Engineering}, QSE) y articula explícitamente seis atributos de
calidad: abstracción, expresividad, verificabilidad, \emph{testing}, escalabilidad y
extensibilidad. La fidelidad de la traducción se instrumenta como un artefacto reproducible
del propio flujo mediante un verificador semántico integrado basado en cálculo ZX, que
compara el grafo derivado de la fuente Quirk con el derivado del OpenQASM generado. Los
resultados muestran que un diseño puede traducirse una única vez a OpenQASM y distribuirse
simultáneamente a varios \emph{backends}, preservando la intención algorítmica del diseño y
habilitando la verificación temprana en la etapa más económica del flujo. QTransLIFIA se
distribuye como biblioteca instalable con una interfaz de programación y una interfaz de
línea de comandos, y como cuaderno Jupyter para exploración y demostración.
\end{abstract}

\noindent\textbf{Palabras clave:} traducción de circuitos cuánticos; Quirk; OpenQASM;
ingeniería de software cuántica (QSE); preservación semántica; compuertas personalizadas;
equivalencia por cálculo ZX; verificación y validación; migración multi-ecosistema; Qiskit;
AWS Braket; PennyLane.

% ===========================================================================
\section{Introducción}
% ===========================================================================

La programación de circuitos cuánticos se desarrolla hoy en un ecosistema
\emph{multi-lenguaje} y esencialmente \emph{multi-vendor}. El diseño conceptual suele
realizarse en herramientas visuales e interactivas como \emph{Quirk}~\cite{gidney2016quirk};
la simulación, compilación y ejecución se delegan a SDK propietarios o de comunidad como
IBM Qiskit~\cite{qiskit}, AWS Braket~\cite{braket} o Xanadu PennyLane~\cite{pennylane}; y
entre ambos extremos circulan lenguajes intermedios de intercambio como
OpenQASM~\cite{cross2017openqasm, openqasm3}. Esta separación de responsabilidades es
deseable desde el punto de vista de la ingeniería, pero introduce un riesgo central: la
\emph{pérdida de fidelidad semántica} durante la traducción. Un circuito que se traduce
incorrectamente ---un control desplazado, un objetivo intercambiado, una abstracción
aplanada o una compuerta descartada--- produce una transformación unitaria distinta de la
diseñada, y el error se propaga silenciosamente a todos los artefactos posteriores.

La \emph{Ingeniería de Software Cuántica} (\emph{Quantum Software Engineering},
QSE)~\cite{qse} propone tratar los artefactos cuánticos con los mismos principios de
calidad que el software clásico: modularidad, trazabilidad, verificabilidad y
reproducibilidad. Bajo esta óptica, un traductor de circuitos no debe limitarse a emitir
sintaxis válida, sino que debe \emph{preservar} la intención algorítmica del diseño y
\emph{demostrar} que lo hace.

\textbf{QTransLIFIA} responde a esta necesidad. Se trata de un traductor que toma circuitos
diseñados en Quirk ---a partir de su URL con el JSON embebido--- y los proyecta sobre
OpenQASM, empleado como lenguaje pivote de representación intermedia. Sobre esa base común,
el sistema genera código ejecutable para varios ecosistemas de destino y verifica la
equivalencia entre el circuito de origen y el traducido mediante cálculo ZX.

\paragraph{Aportes.} Este trabajo presenta:
\begin{enumerate}[label=(\roman*), leftmargin=2em]
  \item un enfoque de \emph{traducción semántica con preservación de abstracciones}, que
        conserva las compuertas personalizadas de Quirk como entidades funcionales
        trasladables, en lugar de aplanarlas;
  \item la instrumentación explícita de seis atributos de calidad de QSE ---abstracción,
        expresividad, verificabilidad, \emph{testing}, escalabilidad y
        extensibilidad--- como decisiones de arquitectura observables;
  \item un \emph{verificador semántico integrado} basado en cálculo ZX~\cite{kissinger2020pyzx},
        que convierte la comprobación de equivalencia en un artefacto reproducible del
        propio flujo;
  \item una arquitectura de representación intermedia que permite migrar un único circuito
        a múltiples \emph{backends} sin reescribir el núcleo de traducción; y
  \item una implementación distribuible como biblioteca (API y CLI) y como cuaderno Jupyter,
        con una batería de pruebas automatizadas sobre circuitos canónicos.
\end{enumerate}

El resto del documento se organiza y estructura de la siguiente manera: la Sección~\ref{sec:arquitectura} describe
Quirk, OpenQASM y la arquitectura del flujo. La Sección~\ref{sec:atributos} desarrolla los
seis atributos de calidad desde la perspectiva QSE. La Sección~\ref{sec:verificacion}
detalla los mecanismos de verificación y \emph{testing}. La
Sección~\ref{sec:migracion} aborda la migración a ecosistemas destino. La
Sección~\ref{sec:discusion} discute implicaciones y limitaciones, y la
Sección~\ref{sec:conclusiones} presenta las conclusiones.

% ===========================================================================
\section{Antecedentes y arquitectura}
\label{sec:arquitectura}
% ===========================================================================

\subsection{El punto de partida: Quirk}

\emph{Quirk}~\cite{gidney2016quirk} es un simulador de circuitos cuánticos basado en
navegador, orientado a la exploración interactiva. Su unidad de diseño es una matriz de
\emph{columnas}; cada columna representa una tajada temporal de operaciones. Un circuito se
almacena íntegramente en la URL de la página como un objeto JSON codificado en el fragmento \code{\#circuit=...}. Quirk no se limita a ofrecer compuertas nativas: permite al usuario \emph{definir sus propias compuertas} de dos maneras:

\begin{itemize}[leftmargin=1.6em]
  \item \textbf{por circuito} (\code{circuit.cols}): una compuerta compuesta por columnas internas que reutilizan otras compuertas (nativas o previamente definidas) que el usuario determina mediante un subrango de segmentos de compuertas seleccionado; y
  \item \textbf{por gates} (\code{matrix}): una compuerta nombrada, especificada directamente por su matriz unitaria y referenciada en el circuito mediante ese nombre.
\end{itemize}

Estas definiciones constituyen el \emph{lenguaje expresivo del diseñador}: la forma en que
el usuario estructura un algoritmo en bloques reutilizables.

\subsection{El lenguaje pivote: OpenQASM}

OpenQASM~\cite{cross2017openqasm, openqasm3} es un lenguaje textual, imperativo y basado en
registros cuánticos (\code{q[...]}), ampliamente soportado por simuladores, transpiladores,
compiladores y SDK. Su adopción como \emph{representación intermedia} convierte a
QTransLIFIA en un puente entre el dominio de diseño visual (Quirk) y los dominios de
ejecución. El sistema soporta dos versiones: OpenQASM~2.0 y OpenQASM~3.0 (esta última por
defecto), seleccionables mediante configuración. La versión~2.0 se obtiene por
transpilación a un conjunto de primitivas equivalente, mientras que la versión~3.0
preserva además la estructura declarativa de las compuertas definidas por el usuario.

\subsection{Arquitectura del flujo}

QTransLIFIA se organiza como una tubería de etapas con responsabilidades independientes, de
modo que la lógica de traducción resulte auditable y reutilizable. La Tabla~\ref{tab:pipeline}
resume las etapas principales y su artefacto característico.

\begin{table}[htbp]
\centering
\small
\caption{Etapas del flujo de traducción y artefacto principal.}
\label{tab:pipeline}
\begin{tabular}{@{}p{0.30\textwidth}p{0.62\textwidth}@{}}
\toprule
\textbf{Etapa} & \textbf{Artefacto / comprobación} \\
\midrule
Análisis del origen & Modelo estructurado del circuito Quirk (dimensiones, compuertas, compuertas personalizadas). \\
Traducción a OpenQASM & Código OpenQASM (\code{algorithms\_qasm/}) como representación intermedia. \\
Verificación semántica & Comparación de grafos ZX (Quirk frente a OpenQASM). \\
Comprobación visual & Capturas y diagramas en PNG, con comparación lado a lado en Markdown. \\
Migración a destino & Scripts Python ejecutables por ecosistema (\code{migrated\_circuits/}). \\
\bottomrule
\end{tabular}
\end{table}

El diseño establece una \emph{frontera clara entre dominios}: el modelo visual y columnar de
Quirk queda contenido en el analizador y el resolutor de compuertas personalizadas,
mientras que la capa de emisión se limita a producir OpenQASM válido. Los detalles de un
dominio no contaminan al otro, en aplicación directa del principio de separación de
preocupaciones.

La arquitectura evita deliberadamente la duplicación de lógica: la traducción
Quirk\,$\rightarrow$\,OpenQASM, la migración a proveedores, la generación de grafos ZX y las
utilidades de carga y dibujo residen en un único módulo de utilidades compartido por las
tres formas de ejecución del sistema ---biblioteca/API, interfaz de línea de comandos y
cuaderno Jupyter---. De este modo, cualquiera de las tres rutas ejecuta exactamente la misma
lógica, sin variantes contradictorias.

\subsection{Formas de ejecución y contrato de entrada}

El flujo funcional es único y se ofrece mediante tres puntos de entrada:

\begin{itemize}[leftmargin=1.6em]
  \item \textbf{Librería (API):} una función de traducción que recibe un \emph{payload} JSON
        y devuelve el resultado sin tocar disco, y una función de ejecución del flujo
        completo sobre el directorio de entrada;
  \item \textbf{Línea de comandos (CLI):} invocable como módulo, con parámetros para la
        entrada, la versión de OpenQASM, los proveedores de destino y los interruptores de
        cada etapa; y
  \item \textbf{Cuaderno Jupyter:} un flujo exploratorio y visual, útil para demostraciones
        y para la inspección artefacto por artefacto.
\end{itemize}

La entrada se especifica como uno o más archivos JSON en los que cada clave es un nombre de
algoritmo y cada valor contiene una URL de Quirk, un desplazamiento de índices opcional y
una descripción opcional. Los artefactos se agrupan por corrida en un directorio con marca
temporal, de forma que las ejecuciones dentro de la misma ventana temporal se consolidan y
las corridas históricas quedan preservadas.

% ===========================================================================
\section{Atributos de calidad desde la ingeniería de software cuántica}
\label{sec:atributos}
% ===========================================================================

QTransLIFIA se concibe con atributos de calidad explícitos, documentados de manera
individual y consolidados como un conjunto coherente. La Tabla~\ref{tab:atributos} los
resume y anticipa cómo se materializa cada uno.

\begin{table}[htbp]
\centering
\small
\caption{Atributos de calidad y su materialización en QTransLIFIA.}
\label{tab:atributos}
\begin{tabular}{@{}p{0.20\textwidth}p{0.72\textwidth}@{}}
\toprule
\textbf{Atributo} & \textbf{Materialización} \\
\midrule
Abstracción & Traducción semántica con preservación de compuertas personalizadas; separación entre análisis y generación; ocultamiento de la complejidad del formato de Quirk. \\
Expresividad & Conservación funcional de las compuertas definibles en Quirk (por circuito y por matriz), respetando el punto de partida de diseño. \\
Verificabilidad & Equivalencia estructural y funcional entre el circuito de Quirk y el OpenQASM generado, apoyada en metadatos y determinismo. \\
\emph{Testing} & Estrategias complementarias por etapa: verificación estática, equivalencia por cálculo ZX y comprobación visual. \\
Escalabilidad & Crecimiento vertical (complejidad del circuito) y horizontal (múltiples ecosistemas) sobre OpenQASM como interfaz estable. \\
Extensibilidad & Incorporación de nuevas compuertas, versiones de OpenQASM y consumidores sin rediseñar el núcleo. \\
\bottomrule
\end{tabular}
\end{table}

\subsection{Abstracción}

La abstracción se manifiesta en tres niveles. Primero, en la \emph{conservación de las
abstracciones del usuario}: el traductor no aplana las compuertas personalizadas, sino que
las mantiene como entidades identificables en OpenQASM. Segundo, en la \emph{ocultación de
la complejidad del formato}: el JSON anidado dentro de la URL de Quirk queda encapsulado por
el analizador, que expone al resto del sistema estructuras de datos normalizadas de alto
nivel (dimensiones del circuito, inventario de compuertas y de compuertas personalizadas).
Tercero, en la \emph{separación entre análisis y generación}: puede analizarse el circuito
sin emitir código, lo que habilita inspección, validación y métricas con independencia de la
sintaxis generada. Esta separación constituye una frontera explícita entre dos modelos
semánticos ---el visual y columnar de Quirk y el textual e imperativo de OpenQASM---
evitando que los detalles de representación de un dominio se filtren en el otro.

\subsection{Expresividad}

La expresividad es la capacidad de \emph{conservar el poder expresivo con el que el usuario
diseñó el circuito} y trasladarlo, sin pérdida, al lenguaje pivote y al ecosistema destino.
El aspecto central es la conservación funcional de las compuertas definibles en Quirk:

\begin{itemize}[leftmargin=1.6em]
  \item \textbf{Compuertas por circuito:} se traducen como declaraciones de compuerta de
        OpenQASM; el cuerpo interno se deriva de sus columnas y se conserva su estructura en
        lugar de expandirse y desaparecer. Las llamadas y las declaraciones coinciden
        siempre, de modo que la compuerta permanece \emph{trasladable y reutilizable} fuera
        de Quirk.
  \item \textbf{Compuertas por matriz:} se convierten en la operación OpenQASM equivalente,
        preservando su comportamiento funcional. Ninguna vía de definición de Quirk queda
        fuera del alcance del traductor.
  \item \textbf{Composición y anidamiento:} las compuertas pueden componerse jerárquicamente
        (una compuerta definida puede contener otras), con un límite de profundidad de
        anidamiento como salvaguarda de robustez. Las instancias multicolumna de una misma
        definición se desambiguan mediante un sufijo que preserva su contexto de uso, y el
        orden de las declaraciones se resuelve topológicamente para respetar la regla
        ``declarar antes de usar''.
\end{itemize}

La conservación se comprueba, no se asume: el modelo estructurado del circuito permite
auditar que toda compuerta definida fue considerada y mapeada, y el cálculo ZX corrobora la
equivalencia semántica de las compuertas trasladadas. El resultado es que los ingenieros
cuánticos pueden definir o importar en Quirk sus compuertas personalizadas ---ya probadas---
sin que la traducción al ecosistema destino constituya una limitante, sino una
\emph{transferencia de módulos funcionales}.

\subsection{Verificabilidad}

La verificabilidad es la capacidad de comprobar que el circuito traducido es equivalente al
original, tanto estructural como funcionalmente. El punto de comprobación más sensible se
ubica en las \emph{etapas iniciales} del flujo: cuando el circuito existe todavía en su
forma de origen. Detectar allí un defecto de traducción es la barrera de control de calidad
más temprana y económica, pues cualquier error introducido en esa fase se propaga a todos
los artefactos posteriores.

La verificación se estructura en tres niveles complementarios:

\begin{description}[leftmargin=1.6em]
  \item[Nivel sintáctico.] El OpenQASM emitido debe ser conforme al estándar: cabecera,
        declaraciones de registros válidas, aridad correcta e identificadores válidos. Se
        resuelve con analizadores OpenQASM existentes.
  \item[Nivel estructural.] Deben coincidir dimensiones, inventario de compuertas,
        controles, objetivos, conectividad y jerarquía de compuertas personalizadas entre el
        circuito analizado y el código emitido. Es totalmente automatizable sin simulación.
  \item[Nivel semántico.] Ambas representaciones deben producir la misma acción unitaria
        sobre el espacio de estados, salvo fase global. Se comprueba de forma nativa mediante
        el cálculo ZX (Sección~\ref{sec:verificacion}).
\end{description}

La preservación estructural es condición \emph{necesaria pero no suficiente} de la
equivalencia funcional: garantiza que no se perdieron piezas, mientras que la comprobación
definitiva de la unitaria la aporta el nivel semántico. La verificabilidad se apoya, además,
en el \emph{determinismo} del traductor ---una misma URL produce siempre el mismo
OpenQASM---, que habilita pruebas de regresión robustas, y en la \emph{trazabilidad} de cada
instrucción hasta la columna y la definición de compuerta que le dieron origen.

\subsection{Testing}

El flujo no es una transformación monolítica, sino una tubería de etapas, y cada
transformación es una oportunidad de introducir un defecto. QTransLIFIA articula tres
familias de estrategias complementarias:

\begin{enumerate}[label=(\arabic*), leftmargin=2em]
  \item \textbf{Verificación estática:} aserción de invariantes sobre el modelo estructurado
        ---número efectivo de qubits, número de columnas, inventario y conteo de compuertas,
        compuertas personalizadas, geometría del \code{offset} y preservación de
        parámetros--- contrastados con el código emitido. Es barata, determinista y
        totalmente automatizable, por lo que se aplica en cada corrida.
  \item \textbf{Equivalencia semántica por cálculo ZX:} comparación automática de la acción
        unitaria de ambas representaciones mediante grafos ZX, que supera la comprobación
        estructural al detectar estados de qubits distintos, controles duplicados o
        expansiones incorrectas de compuertas personalizadas.
  \item \textbf{Comprobación visual rápida:} capturas del circuito de Quirk, diagramas del
        OpenQASM generado y comparaciones lado a lado, aptas para inspección humana
        inmediata en circuitos sencillos y como \emph{sanity check} en circuitos grandes.
\end{enumerate}

Ninguna estrategia es suficiente por sí sola: la estática garantiza que no se pierden
piezas, la semántica aporta la garantía formal de equivalencia y la visual ofrece la lectura
inmediata que acelera la detección de defectos. El repositorio complementa estas
estrategias con una suite de pruebas automatizadas que ejercita los modos de entrada, los
circuitos canónicos, los desplazamientos de índice, las compuertas personalizadas y la
opción de preservación de compuertas nativas.

\subsection{Escalabilidad}

La escalabilidad se aborda en dos direcciones. La \emph{vertical} corresponde al crecimiento
en la complejidad del circuito: más qubits, más columnas y jerarquías de compuertas
personalizadas más profundas. Se soporta mediante un cálculo de alturas consciente de la
geometría ---que determina la altura ocupada por cada compuerta multi-qubit o compuesta y
evita asignaciones incorrectas de índices--- y mediante un cálculo dinámico del número
efectivo de qubits y del desplazamiento, de modo que la asignación a registros permanezca
coherente aunque las abstracciones añadan qubits virtuales o modifiquen la altura efectiva.

La escalabilidad \emph{horizontal} corresponde a la integración con múltiples ecosistemas
destino. El eje central es OpenQASM como interfaz estable: al emitir este formato
intermedio, el traductor queda agnóstico al \emph{backend}. Nuevos consumidores ---
simuladores, transpiladores, \emph{toolchains} de validación--- pueden integrarse como
consumidores puros del OpenQASM generado, y es factible añadir etapas posteriores
(optimización, mapeo a topología física, reescritura para dispositivos específicos) que
operen sobre el IR sin acoplarse a la lógica de análisis de Quirk.

\subsection{Extensibilidad}

La extensibilidad es la capacidad de evolucionar sin rediseñar la arquitectura central.
Se fundamenta en tres principios: separación de responsabilidades ---análisis, expansión de
compuertas personalizadas y emisión están claramente desacoplados---; abstracción mediante
mapeos, donde una tabla de compuertas y un modelo estructurado del circuito facilitan añadir
elementos sin alterar los mecanismos existentes; y centralización de la emisión, que permite
sustituir o extender la capa de salida de forma aislada.

Estos principios habilitan cuatro líneas de evolución:

\begin{itemize}[leftmargin=1.6em]
  \item \textbf{Evolución de Quirk:} el mapeo por identificador y nombre absorbe cambios de
        nomenclatura o nuevos identificadores, y la extracción explícita de compuertas
        personalizadas hace que el soporte a nuevas abstracciones sea implícito.
  \item \textbf{Nuevas compuertas y funcionalidades:} incorporar compuertas nativas localiza
        el cambio en la capa de mapeo y emisión, sin afectar el cálculo de alturas, el
        mecanismo de expansión ni el análisis. La información estructurada del circuito
        habilita, además, funcionalidades futuras ---validación, métricas, reportes
        estructurados--- sin acoplarlas a la generación.
  \item \textbf{Nuevas versiones de OpenQASM:} la emisión aislada permite introducir un nuevo
        generador reutilizando íntegramente el análisis y la expansión de abstracciones. La
        coexistencia de OpenQASM~2.0 y 3.0 es una manifestación temprana de esta capacidad.
  \item \textbf{Nuevos ecosistemas destino:} al adoptar un IR canónico, los nuevos destinos
        se integran como adaptadores o consumidores externos, sin proliferar lógica específica
        de Quirk por \emph{backend}.
\end{itemize}

% ===========================================================================
\section{Verificación y validación}
\label{sec:verificacion}
% ===========================================================================

\subsection{El verificador semántico integrado}

El componente más distintivo de QTransLIFIA es su verificador semántico basado en cálculo
ZX, que se apoya en la librería PyZX~\cite{kissinger2020pyzx}. El proceso consta de tres
pasos: se construye un grafo ZX \emph{directamente desde la fuente Quirk}; se construye un
segundo grafo \emph{desde el OpenQASM generado}; y se comparan ambos grafos por circuito. La
estrategia de comparación se adapta al tamaño del circuito: para circuitos pequeños se
reduce cada grafo y se comparan sus tensores; para circuitos grandes se verifica la igualdad
de los circuitos asociados considerando la fase global. El resultado por circuito
---\emph{equivalente}, \emph{no equivalente} o \emph{no comprobado}--- se registra como un
reporte auditable.

La regla de equivalencia y el criterio de decisión residen en un único punto del sistema,
compartido por las tres formas de ejecución, de modo que el veredicto es idéntico
independientemente de la ruta empleada. Además, el verificador se protege con un límite de
tiempo por comparación: cuando la reducción de un grafo resulta computacionalmente
intratable, el circuito se marca como \emph{no comprobado} en lugar de bloquear el flujo.
Este mecanismo convierte la comprobación semántica ---el nivel más fuerte de verificación---
en un \emph{artefacto reproducible del propio pipeline}, en línea con la exigencia de la QSE
de instrumentar la verificación como parte del proceso de desarrollo y no como una actividad
manual posterior.

\subsection{Estrategias automatizables}

La separación entre análisis y generación habilita estrategias de verificación automatizada
adicionales: el contraste entre el inventario de compuertas del circuito analizado y las
instrucciones efectivamente emitidas; la comprobación de invariantes (\emph{property-based
testing}) sobre dimensiones, número de controles y rango de índices; las pruebas
metamórficas, que exigen invariancia ante transformaciones que no alteran el circuito; y la
validación sintáctica con consumidores externos. El determinismo del traductor sustituye,
en estas estrategias, la necesidad de comparar contra una referencia móvil: la comparación
contra una salida fija es estable a lo largo del tiempo.

\subsection{Estado de la validación}

La batería de pruebas del repositorio ejercita los modos de entrada individual y por lotes,
circuitos canónicos (estados de Bell, compuertas de Toffoli, compuertas personalizadas
simples y anidadas), el desplazamiento de índices y la opción de preservación o expansión de
compuertas. La verificación semántica por cálculo ZX se ejecuta sobre las colecciones de
algoritmos de prueba, que abarcan algoritmos principales, variantes populares y diseños
específicos de la transformada cuántica de Fourier y su inversa. El flujo completo ---
incluidas la generación de imágenes, la comparación visual, los grafos ZX y la migración a
proveedores--- se ejecuta de extremo a extremo tanto desde la línea de comandos como desde el
cuaderno, sin intervención manual.

% ===========================================================================
\section{Migración a ecosistemas destino}
\label{sec:migracion}
% ===========================================================================

Sobre la representación intermedia OpenQASM, QTransLIFIA genera scripts Python ejecutables,
uno por algoritmo y ecosistema. La Tabla~\ref{tab:proveedores} resume los destinos
soportados y sus particularidades.

\begin{table}[htbp]
\centering
\small
\caption{Ecosistemas de destino soportados.}
\label{tab:proveedores}
\begin{tabular}{@{}p{0.20\textwidth}p{0.68\textwidth}@{}}
\toprule
\textbf{Ecosistema} & \textbf{Particularidad de la migración} \\
\midrule
IBM Qiskit & Emplea compuertas nativas siempre que la descomposición del circuito corresponda
a una puerta estándar; produce además diagramas con el estilo de IBM Quantum Composer. \\
AWS Braket & Requiere el SDK de Amazon Braket; el cargador nativo se complementa con un
respaldo por migración para cubrir todo el OpenQASM que genera el proyecto. \\
Xanadu PennyLane & Se materializa siempre a través del script migrado, generando circuitos
ejecutables con un número configurable de disparos. \\
\bottomrule
\end{tabular}
\end{table}

La migración se beneficia de la preservación de abstracciones: las compuertas definidas por
el usuario llegan al SDK destino con su semántica intacta ---recreadas como compuertas
personalizadas--- o reducidas a primitivas equivalentes por el propio SDK. Un mismo circuito
puede migrarse una única vez a OpenQASM y distribuirse a varios ecosistemas
simultáneamente, de modo que la expresividad conservada rinde en todos los destinos a la
vez. La carga del circuito OpenQASM por proveedor también reside en un único punto
compartido, de forma que las tres formas de ejecución interpretan el IR de manera idéntica.

% ===========================================================================
\section{Discusión}
\label{sec:discusion}
% ===========================================================================

El aporte conceptual de QTransLIFIA no está en la conversión de sintaxis ---resoluble con
mapeos locales--- sino en sostener dos premisas de ingeniería. La primera es que
\emph{la abstracción del usuario es parte de la semántica}: aplanar una compuerta
personalizada durante la traducción destruye información de diseño y dificulta el
mantenimiento del software cuántico. La segunda es que \emph{la fidelidad debe ser
demostrable}: incorporar el verificador semántico al flujo convierte una propiedad deseable en
una evidencia reproducible.

Entre las limitaciones que delinean el trabajo futuro se destacan: la verificación ZX puede
resultar intratable en circuitos con muchas compuertas multi-controladas, motivo por el cual
se emplea un límite de tiempo y un veredicto explícito de \emph{no comprobado}; la
equivalencia se establece salvo fase global, lo que es adecuado para la mayoría de los
propósitos pero conviene documentar por contexto; y la migración a nuevos ecosistemas, si
bien no exige modificar el núcleo, sí requiere implementar un generador específico por
destino. Finalmente, la captura visual del circuito original depende de un navegador
automatizado, por lo que es la etapa más sensible al entorno de ejecución.

Como líneas de evolución, la arquitectura admite incorporar nuevos ecosistemas como
consumidores puros del IR, añadir nuevos generadores para versiones futuras de OpenQASM sin
rediseñar el análisis, y enriquecer el modelo estructurado con nuevas métricas de calidad.

% ===========================================================================
\section{Conclusiones}
\label{sec:conclusiones}
% ===========================================================================

QTransLIFIA demuestra que es posible migrar circuitos cuánticos entre plataformas
preservando la \emph{intención algorítmica} del diseño y \emph{demostrando} dicha
preservación como parte del propio flujo. Al combinar traducción semántica con preservación
de abstracciones, OpenQASM como representación intermedia estable y un verificador
semántico integrado basado en cálculo ZX, el sistema materializa los seis atributos de
calidad que la Ingeniería de Software Cuántica promueve: \emph{expresividad} y
\emph{abstracción} definen qué se conserva; \emph{verificabilidad} y \emph{testing}
garantizan cómo se comprueba; y \emph{escalabilidad} y \emph{extensibilidad} delimitan hasta
dónde puede crecer.

Tres resultados sostienen esta conclusión. Primero, las compuertas personalizadas definibles
en Quirk ---por circuito o por matriz--- sobreviven a la traducción como entidades
funcionales, reutilizables y auditables, en lugar de aplanarse. Segundo, la equivalencia
entre el circuito diseñado y el traducido se convierte en un artefacto reproducible del
pipeline, no en una comprobación manual posterior. Tercero, un único circuito puede
migrarse a varios ecosistemas sin reescribir el núcleo, gracias a una arquitectura modular
desacoplada y a un lenguaje pivote ampliamente soportado.

Estos tres resultados convergen en la premisa central del proyecto:
\textbf{el circuito que el usuario diseñó en Quirk es el circuito que se traduce, se verifica
y se ejecuta} en cualquiera de los ecosistemas de destino. Al sentar esta premisa sobre
principios explícitos de ingeniería de software cuántica, QTransLIFIA se posiciona como una
solución preparada para evolucionar con el ecosistema cuántico, preservando al mismo tiempo
estabilidad, mantenibilidad y fidelidad semántica.

% ===========================================================================
\begin{thebibliography}{9}
% ===========================================================================

\bibitem{gidney2016quirk}
C.~Gidney, \emph{Quirk: a drag-and-drop quantum circuit simulator}.
\url{https://algassert.com/quirk}. Accedido en 2026.

\bibitem{cross2017openqasm}
A.~W. Cross, L.~S. Bishop, J.~A. Smolin y J.~M. Gambetta,
\emph{Open Quantum Assembly Language}, 2017. \url{https://openqasm.com}.

\bibitem{openqasm3}
A.~W. Cross \emph{et al.}, \emph{OpenQASM 3: A broader and deeper quantum assembly language},
ACM Transactions on Quantum Computing, 2022.

\bibitem{kissinger2020pyzx}
A.~Kissinger y J.~van de Wetering, \emph{PyZX: A Python library for quantum circuit
rewriting and simplification}, 2020. \url{https://pyzx.readthedocs.io}.

\bibitem{qiskit}
Qiskit Development Team, \emph{Qiskit: an open-source framework for quantum computing}.
\url{https://www.ibm.com/quantum/qiskit}. Accedido en 2026.

\bibitem{braket}
Amazon Web Services, \emph{Amazon Braket: servicio de computación cuántica gestionada}.
\url{https://aws.amazon.com/braket/}. Accedido en 2026.

\bibitem{pennylane}
Xanadu, \emph{PennyLane: cross-platform Python library for quantum machine learning and
quantum computing}. \url{https://pennylane.ai}. Accedido en 2026.

\bibitem{qse}
M.~Piattini, G.~Peterssen y R.~Pérez-Castillo, \emph{Quantum Software Engineering},
Springer, 2022.

\end{thebibliography}

\end{document}
